\documentclass[12pt]{article}

\usepackage[a4paper,margin=1.5cm]{geometry}
\usepackage{booktabs}
\usepackage{hyperref}

\begin{document}
	\title{\LARGE{HUMAN VALUE AND PROFESSIONAL ETHICS PYQ}}
	\maketitle
	
	
	\newpage
	
	\section{\LARGE{PYQS}}
	
	\begin{enumerate}
		\subsection{1-marks questions}
		
		
		
		
		
		
		\subsection{2-marks questions}
		
		\item Discuss the following proposal "I will be bored of happiness if I am always happy". \hfill{[CAT 25]}
		
		\item Explain the term "Respect". \hfill{[CAT 25,END25]}
		
		\item What is importance of the term "proposal" in value education. \hfill{[END 25]}
		
		\item \textit{Self-Exploration} is a process of a dialogue between \textit{"what you are"} and \textit{"what you really want to be"}.Comment on the above line. \hfill{[END 25]}
		
		\item Write the all dimension of human living. \hfill{[END 25]}
		
		\item State any two activites which involve both I \textit{(self)} \& \textit{(body)}. \hfill{[END 25]}
		
		\item Give a list of basic guidelines for value education. \hfill{[CAT 24]}
		
		\item Name all four levels of living. \hfil{[CAT 24]}

		
		\item with help of a neat diagram discuss "The program to fulfill Basic Human Aspiration". \hfill{[CAT 24]}
		
		\item Write one line defination of \textit{respect} and \textit{love}. \hfill{[CAT 24]}
		
		\item Explain the term \textit{"All encompassing"} as basic guidelines for the value education. \hfill{[END 24]}
		
		\item \textit{Self Exploration} is a process of \textit{self-evolution} through \textit{self-investigation}. Comment on above line. \hfill{[END 24]}
		
		\item Write the four levels of living. \hfill{[END 24]}
		
		\item State any two activites which involves only \textit{I (Self)}. \hfill{[END 24]}
		
		\item Define the term \textit{Trust}. \hfill{[END 24]}
		
		\item Mention four levels of living. \hfill{[CAT 22]}
		
		\item Justify that need of \textit{"I"} is qualitative whereas need of body is quantitative. \hfill{[END 22]}
		
		\item What is difference between the terms \textit{Rational \& Universal} as basis for value education? \hfill{[END 22]}
		
		\item "Self-exploration is a process of knowing oneself and throught that,knowing the entire existance."Give an example to justify the above line. \hfill{[END 22]}
		
		\item Write the four level of living. \hfill{[END 22]}
		
		\item State any three activities which involves both "I" and "Body". \hfill{[END 22]}
		
		\item Beside the physical assets,what are the other needs of human Beings. \hfill{[END 22]}
		
		
		
		\subsection{3-marks questions}
		
		\item List all dimension and all levels of human living. \hfill{[CAT 25]}
	
		\item Comment on the statement "Self exploration is a process self evolution through seld investigation". \hfill{[CAT 25]}
		
		\item Name any two activities involving self and body both.Explain how? \hfill{[CAT 25]}
		
		\item List the "Purpose of Self Exploration". \hfill{[CAT 24]}
		
		\item How values differ from skills?Are they complimentary? \hfill{[CAT 22]}
		
		\item What is basic aspiration in following statements? \hfill{[CAT 22]}
		\begin{enumerate}
			\item \textit{"I want to do research in astronomy"}
			\item \textit{"I want to be a film star"}
			\item \textit{"i want to become a doctor"}
		\end{enumerate}
		
		\item Mention any one activity based on \hfill{[CAT 22]}
		\begin{enumerate}
			\item Self
			\item Self and Body
			\item Body with consent self
		\end{enumerate}
		
		\item What do you mean by living in Self-regulation(Sanyama) for achieving harmony in Self (I) and body.
		
		
		\subsection{4-marks questions}

		\item Differentiate needs of I and body with the help of a table. \hfill{[CAT 24]}
		
		\item List the different power and activities of self. \hfill{[CAT 24]}
		
		\item "Physical facilites are necessary and complete for animals,while they are necessary but not complete for humans." Comment. \hfill{[CAT 22]}
		
		
		
		
		
		
		
		
		\subsection{5-marks questions}
		
		\item Write a short note on SANYAMA and SVASTHYA. \hfill{[CAT 25]}
		
		
		
		
		
		
		
		
		
		\subsection{10-marks questions}
		
		\item Discuss the basic guidelines for the value education. \hfill{[END 25]}
		
		\item Wheather the following statement are true or false?Justify. \hfill{[END 25]}
		\begin{enumerate}
			\item \textit{I will be bored of happiness if I am always happy}.
			
			\item \textit{I need to be unhappy to recognize that I am happy.}
		\end{enumerate}
		
		\item Discuss the following power and activities. \hfill{[END 25]}
		\begin{enumerate}
			\item Power
			\begin{enumerate}
				\item \textit{Desire (Ichchah)}
				\item \textit{Thought (Vichara)}
			\end{enumerate}
			
			\item Activity
			\begin{enumerate}
				\item \textit{imaging (Chintrana)}
				\item \textit{Analyzing (Vislesana)}
			\end{enumerate}
		\end{enumerate}
		
		\item Under the holistic criteria for evolution,state the specific criteria for jufging the appropriateness of the technologies and production systems. \hfill{[END 25]}
		
		\item Discuss the comprehensive human goals. \hfill{[END 25]}
		
		
		\item Compare all four orders of the nature \textit{i.e. Material Order, Plants and Bio order,Animal order and Human order.} \hfill{[END 24]}
		
		\item Discuss the specific criteria for judging the appropriateness of production systems. \hfill{[END 24]}
		
		\item What do you mean by term \textit{"Holistic Education"} and \textit{"Humanistic Constitution"}? \hfill{[END 24]}
		
		\item Discuss unit and space and its coexistence. \hfill{[END 24]}
		
		\item Explain the definitiveness of ethical human conduct in terms of values,policies and character. \hfill{[END 24]}
		
		
		\item What is basic human aspiration?Justify your answer. \hfill{[END 22]}
		
		\item differentiate between "I" and "Body" with the help a table. \hfill{[END 22]}
		
		\item Define the following: \hfill{[END 22]}
		\begin{enumerate}
			\item Guidance
			\item Reverence
			\item Care
			\item Glory
			\item Gratitude
		\end{enumerate}
		
		\item What do you mean by "Holistic Education" and "Humanistic Constitution"?
		
		\item What are the comprehensive goals?

		\section{20-marks questions}
		
		\item What is value education?Explain the need to study value education.State the process of the value education. \hfill{[END 25]}
		
		\item Compare the I (Self) and Body on the basis of the following: \hfill{[END 25]}
		\begin{enumerate}
			\item Needs
			\item Activites
		\end{enumerate}
		
		\item What is \textit{Self-exploration}?With a neat diagram explain the program for \textit{self-exploration}. \hfill{[END 24]}
		
		\item Explain the utility of following programme to achieve Comprehensive human goal. \hfill{[END 24]}
		\begin{enumerate}
			\item Education - right Living (Sikshā - Sanskāra)
			\item Justive - Preservation (Nyāya-Surakshā)
			\item Health - Self-regulation (Svāsthya-Sanyama)
			\item Production - Work (Utpādana-Kārya)
			\item Exchange - Storage (Vinimaya - Kosa)
		\end{enumerate}
		
		\item Discuss the process of self-exploration with the help of diagram. \hfill{[END 22]}
		
		\item Discuss five dimension of Human endeavour for ensuring harmony in society. \hfill{[END 22]}
	\end{enumerate}
\end{document}